\section{Size arithmetic of the symmetric \texorpdfstring{$k=7$}{k=7} reduction}\label{sa:sec}

This section gives written proofs for three small Lean files of the research records that no other section covers.
Two of them check the size arithmetic of the reduction of the symmetric $k=7$ statement to its case classes; the third
instantiates Baek's theorem at $7$-gons.

\paragraph{The three files and their role.}
\begin{enumerate}[label=(\arabic*)]
\item \dirn{AUDIT3/A1_reduction/SizeArith.lean} (namespace \lean{A1}, $20$ theorems). The directory \dirn{AUDIT3}
  records an adversarial audit, run on 1~November 2026 against commit \dirn{9cf4949a5}, of the statement ``every set
  of at least $33$ points in general position with a nontrivial symmetry, isometric or affine, contains a convex
  $7$-gon''; it records an audit and claims no theorem (\dirn{AUDIT3/README.md}). Its component audit~A1 (report
  \dirn{AUDIT3/reports/A1_reduction.md}) re-derived by hand the reduction of this statement to the case theorems of
  the records for half-turns, threefold and fivefold rotations and reflections, to which the case classes of
  Definition~\ref{sr:def:classes} correspond. The auditor wrote this file to check in Lean, for all sizes $n$, the
  arithmetic of that reduction and the arithmetic behind the upper bounds of the table of the largest symmetric
  $7$-gon-free sizes (trust-table rows 4, 8 and~9 of the report).
\item \dirn{AUDIT3/A1_reduction/verify_A1-F1/AssembledSizes.lean} ($7$ theorems and one anonymous example). Finding
  A1-F1 of that audit (rated minor) is that the assembled theorem was not written anywhere. The verifier of the
  finding wrote this file to check, for all $n$, the bookkeeping of the assembled statement: the half-turn, threefold
  and fivefold cases, the two mirror cases and the fourfold case through the half-turn, together with a control
  showing that the half-turn bound fails at $32$ points (\dirn{AUDIT3/reports/A1_reduction.md}, lines 103--106). The
  file also states, at its line~24, that an order of at least $7$ gives at least one orbit other than the centre
  (\lean{big_order_orbit}); the report's description of the file, at its line~106, does not mention this lemma.
\item \dirn{ES7SPLIT/review/equivalence/BaekSeven.lean} (one theorem). It instantiates the vendored formalization
  \lean{main} of Baek's theorem, cited in this document as Theorem~\ref{mb:thm:baek}, at $7$-gons: $17$ elements of
  an abstract configuration contain a $4$-cap or a $7$-gon. The records of \dirn{ES7SPLIT} use it as the Lean-checked
  abstract form of the one published theorem that Theorems~5 and~6 and Corollary~7 of \dirn{ES7SPLIT/README.md}
  (written proofs there) use as a black box. It plays no role in the reduction.
\end{enumerate}

\paragraph{What the files do not contain.}
The statements of the first two files are facts about natural numbers. The geometric modelling that connects them
with point sets is not part of them: the header comment of \dirn{SizeArith.lean} (lines 4--8) states it and calls it
``audited by hand in the report, not formalised here''. In that modelling, a rotation of order $s$ about $c$ gives
$n=s\cdot k+e$, where $k$ is the number of orbits other than $\{c\}$, $e=1$ if $c\in S$ and $e=0$ otherwise, so
$e\le1$, and $e=0$ for even $s$; a reflection gives $n=2k+b$, where $k$ is the number of mirror pairs and $b\le2$ the
number of points on the axis; and deleting whole orbits, the centre or axis points keeps the symmetry and general
position. We prove these facts as a written step (Lemma~\ref{sa:lem:model}), which is not a declaration of these
files. Related facts are Lean-checked elsewhere in this document, in other forms: Theorem~\ref{ec:thm:compat}(a),(b),
an abstract form of the count $n=s\cdot k+e$ for a map $z\mapsto\zeta z$ of exact order $s$ in a commutative group
with zero; the selection of whole orbits (Lemma~\ref{sr:lem:orbit}) and the counting of half-turn pairs, rotation
orbits and mirror pairs in the Lean proof of Theorem~\ref{sr:thm:main}; the layer results
(Theorems~\ref{ll:thm:layers} and~\ref{ll:thm:card-peel}), through verbatim copies of which, and of
Theorem~\ref{ll:thm:reduction}, the Lean proofs of Theorems~\ref{ec:thm:fourfold} and~\ref{ec:thm:fivefold} pass
from a symmetric set to a subset of exactly $32$, respectively $35$, points (the file of these proofs repeats
lines 1--1147 of \dirn{LAYERLEAN/LayerLemma.lean} instead of importing it); and Lemmas~\ref{ll:lem:minperiod}
and~\ref{lg:lem:centre}, which the proof of
Lemma~\ref{sa:lem:model} uses. No Lean theorem connects these facts with the statements of the files of this
section. Remark~\ref{sa:rem:use} explains how Lemma~\ref{sa:lem:model} combines with the arithmetic and with the case
classes. The Lean-checked reduction of the symmetric statement to the case classes is Theorem~\ref{sr:thm:main} for
Euclidean isometries, with Theorem~\ref{ar:thm:main} for affine symmetries. Both are conditional on the classes,
and their Lean proofs do their own arithmetic: a text search of the repository's Lean files finds the names of the
declarations of \dirn{SizeArith.lean} and \dirn{AssembledSizes.lean} only in these two files, so no other Lean file
uses them.

\paragraph{Sources and evidence.}
Paths are relative to the research-records directory \dirn{cloud/es7-2026-10-19/results}, except those beginning
with \dirn{ErdosSzekeres/}, which are relative to the repository root. The table lists the three Lean files and
their logs by file name; their full paths are given in the list that follows it. The files are those of commit
\dirn{7e32ee259}; their line counts and SHA-256 hashes (computed with \mbox{\texttt{shasum -a 256}}) are:
\begin{center}\small
\begin{tabular}{@{}lrl@{}}
\toprule
file & lines & SHA-256\\
\midrule
\dirn{SizeArith.lean} & $103$ & \texttt{170384d81db76bb5e8d019e886f50b2545dc7581599f077c6431e08edb532418}\\
\dirn{SizeArith.lean.log} & $20$ & \texttt{b4bca1cd8542a5e39d2a8560570b573ece40533d9850e8b77f5fc33c41da5770}\\
\dirn{AssembledSizes.lean} & $39$ & \texttt{c08bdf200fb59f8353c999460a1069ecd7e3414874c4a79a8030c58dc36f2feb}\\
\dirn{AssembledSizes.lean.log} & $11$ & \texttt{1db4174c7c04edd5a20fad31800eaaf076a0af5cb480c08dcc1771018ae377c9}\\
\dirn{BaekSeven.lean} & $12$ & \texttt{d6d6636c78dd1a2b271ef8cd1dd09a9c869d5db8f3319f702921319eaf36c2d0}\\
\dirn{BaekSeven.out} & $3$ & \texttt{53dced1cd76f1ff0b2afb9ba3ba854ffe0ac1a1dbe895855cb6abfa108f8c747}\\
\bottomrule
\end{tabular}
\end{center}
\begin{itemize}
\item \dirn{AUDIT3/A1_reduction/SizeArith.lean}: namespace \lean{A1}, no \texttt{import} (core Lean only). Lines 1--9
  are the header comment; lines 16--80 hold the $20$ theorems; lines 84--103 are one \texttt{\#print axioms} command
  for each of them, in the same order. Its log \dirn{AUDIT3/A1_reduction/SizeArith.lean.log} is the output of these
  commands and contains no error and no warning. It reports the axioms \lean{propext}, \lean{Classical.choice},
  \lean{Quot.sound} for \lean{A1.c5_reduce}, \lean{A1.c4_reduce}, \lean{A1.c6_reduce}, \lean{A1.c2_reduce},
  \lean{A1.c5_gap} and \lean{A1.table_mirror}, and \lean{propext}, \lean{Quot.sound} for the other $14$ theorems.
\item \dirn{AUDIT3/A1_reduction/verify_A1-F1/AssembledSizes.lean}: root namespace, no \texttt{import}. The seven
  theorems are declared at lines 6, 9, 12, 15, 18, 21 and~24 (the last one with its proof on lines 25--27), the
  anonymous \texttt{example} at line~30, and lines 33--39 are a \texttt{\#print axioms} command for each theorem.
  Its log \dirn{AUDIT3/A1_reduction/verify_A1-F1/AssembledSizes.lean.log} contains the linter warning ``unused
  variable \texttt{hp}'' at line~24, column~41, then \lean{propext}, \lean{Classical.choice},
  \lean{Quot.sound} for \lean{halfturn_pairs} and \lean{propext}, \lean{Quot.sound} for the other six theorems, and a
  timing line; it contains no error.
\item \dirn{ES7SPLIT/review/equivalence/BaekSeven.lean}: imports \lean{ErdosSzekeres.Vendor.ETV.Main.Main} (hence
  Mathlib). Line~6 is \lean{baek_seven}, line~11 the command \mbox{\texttt{\#check @main}}, line~12 the command
  \mbox{\texttt{\#print axioms baek\_seven}}. Its output \dirn{ES7SPLIT/review/equivalence/BaekSeven.out} gives in
  lines 1--2 the type of \lean{main} (Remark~\ref{sa:rem:baek-seven}) and in line~3 the axioms \lean{propext},
  \lean{Classical.choice}, \lean{Quot.sound} of \lean{baek_seven}. Since \texttt{\#print axioms} reports every axiom
  in the dependency closure of a declaration, this line also covers the vendored theorem \lean{main}
  (\dirn{ErdosSzekeres/Vendor/ETV/Main/Main.lean}, SHA-256 prefix \texttt{83dd0a214204}) as it is used.
\end{itemize}
The first two files were checked with Lean~4.24.0 alone (\dirn{AUDIT3/README.md}, section ``Reproduction'';
\dirn{AUDIT3/reports/A1_reduction.md}, lines 106 and 262), the third with \mbox{\texttt{lake env lean}} from the
repository root (\dirn{ES7SPLIT/review/README.md}, section ``Reproduction''), that is, with the repository's
Lean~4.24.0 and
Mathlib \dirn{f897ebcf}. A text search of the three files finds no \texttt{sorry}, \texttt{admit},
\texttt{native\_decide} or \texttt{axiom} declaration. The logs were produced by the project and were not re-run for
this document.

\paragraph{Rigor.}
Every statement below whose head names a Lean declaration is \emph{Lean-checked}, with the axioms shown in the logs
above (all within \lean{propext}, \lean{Classical.choice}, \lean{Quot.sound}). The proofs given here are
\emph{written proofs (not refereed)}. No statement of this section is computer-assisted. None of them is a finite
computation: the Lean-checked statements are universally quantified, over natural numbers or, in
Theorem~\ref{sa:thm:baek-seven}, over linearly ordered types, configurations and finite sets, and the control of
Remark~\ref{sa:rem:control} is the negation of such a statement. The Lean proofs use \texttt{decide} only to close
side facts inside these proofs, never to evaluate a whole statement: the five
closed instances $s=2,\dots,6$ of the conclusion of Lemma~\ref{sa:lem:split}, the value $\binom62+2=17$ in
Theorem~\ref{sa:thm:baek-seven} and $32\le32$ in the control; the written proofs check these facts by hand as well.
By the convention that accompanies the rigor labels at the beginning of this document, such a use of \texttt{decide}
does not make a statement computer-assisted.
Lemma~\ref{sa:lem:model} and Remark~\ref{sa:rem:use} are written steps and are not Lean-checked.
Theorem~\ref{sa:thm:baek-seven} rests on Baek's theorem, Theorem~\ref{mb:thm:baek}, which is cited and
not proved in this document.

\paragraph{Conventions.}
In Subsections~\ref{sa:ssec:lower}--\ref{sa:ssec:assembled} all variables range over $\N=\{0,1,2,\dots\}$ (Lean
\texttt{Nat}), and each statement is universally quantified over its variables in the order in which Lean lists
them. No statement subtracts natural numbers. Lean's product \mbox{\texttt{2 * t * k}} is $(2t)k$ (multiplication
associates to the left). The Lean proofs use the tactic \texttt{omega}, a decision procedure for linear arithmetic
over $\N$ and
$\Z$ whose proof term is checked by the kernel; the written proofs are direct arguments.

%----------------------------------------------------------------------------------------------------------------------
\subsection{Lower bounds for the reduction}\label{sa:ssec:lower}

The first eleven theorems of \dirn{SizeArith.lean} (its part ``Reduction'', lines 13--56) turn ``at least $33$
points'' into the numbers of orbits or pairs that the case classes need (Remark~\ref{sa:rem:use}).

\begin{proposition}[rotations of order $2$ to $6$;\\
  \leanref{A1.c3_reduce}{AUDIT3/A1_reduction/SizeArith.lean:16},\\
  \leanref{A1.c5_reduce}{AUDIT3/A1_reduction/SizeArith.lean:19},\\
  \leanref{A1.c4_reduce}{AUDIT3/A1_reduction/SizeArith.lean:22},\\
  \leanref{A1.c6_reduce}{AUDIT3/A1_reduction/SizeArith.lean:25},\\
  \leanref{A1.c2_reduce}{AUDIT3/A1_reduction/SizeArith.lean:28}]\label{sa:prop:rot}
All variables range over $\N$.
\begin{enumerate}[label=(\alph*)]
\item \textup{(}\lean{A1.c3_reduce}\textup{)} For all $k,e\in\N$: if $e\le1$ and $33\le3k+e$, then $11\le k$.
\item \textup{(}\lean{A1.c5_reduce}\textup{)} For all $k,e\in\N$: if $e\le1$ and $33\le5k+e$, then $35\le5k+e$ and
  $7\le k$.
\item \textup{(}\lean{A1.c4_reduce}\textup{)} For every $k\in\N$: if $33\le4k$, then $36\le4k$, $8\le k$ and
  $17\le2k$.
\item \textup{(}\lean{A1.c6_reduce}\textup{)} For every $k\in\N$: if $33\le6k$, then $11\le2k$ and $17\le3k$.
\item \textup{(}\lean{A1.c2_reduce}\textup{)} For every $k\in\N$: if $33\le2k$, then $34\le2k$ and $17\le k$.
\end{enumerate}
\end{proposition}

\begin{proof}
(a) From $33\le3k+e$ and $e\le1$ we get $3k\ge32$; since $3\cdot10=30<32$, $k\ge11$.
(b) In the same way $5k\ge32$, and $5\cdot6=30<32$ gives $k\ge7$; then $5k+e\ge5k\ge35$.
(c) Since $4\cdot8=32<33\le4k$, we have $k\ge9$; hence $4k\ge36$, $k\ge8$ and $2k\ge18\ge17$. (So the statement records
the weaker bound $8\le k$.)
(d) Since $6\cdot5=30<33\le6k$, we have $k\ge6$; hence $2k\ge12\ge11$ and $3k\ge18\ge17$.
(e) Since $2\cdot16=32<33\le2k$, we have $k\ge17$, and then $2k\ge34$.
Lean proves each part with \texttt{omega}.
\end{proof}

\begin{proposition}[reflections;\\
  \leanref{A1.mirror_reduce_b0}{AUDIT3/A1_reduction/SizeArith.lean:31},\\
  \leanref{A1.mirror_reduce_b12}{AUDIT3/A1_reduction/SizeArith.lean:34}]\label{sa:prop:mirror}
All variables range over $\N$.
\begin{enumerate}[label=(\alph*)]
\item \textup{(}\lean{A1.mirror_reduce_b0}\textup{)} For every $k\in\N$: if $33\le2k$, then $17\le k$.
\item \textup{(}\lean{A1.mirror_reduce_b12}\textup{)} For all $k,b\in\N$: if $1\le b$, $b\le2$ and $33\le2k+b$,
  then $16\le k$.
\end{enumerate}
\end{proposition}

\begin{proof}
(a) As in Proposition~\ref{sa:prop:rot}(e), $2\cdot16=32<33\le2k$ gives $k\ge17$.
(b) From $33\le2k+b$ and $b\le2$ we get $2k\ge31$; since $2\cdot15=30<31$, $k\ge16$. The hypothesis $1\le b$ is not
used; it records that the statement is meant for one or two points on the axis. Lean proves both parts with
\texttt{omega}.
\end{proof}

\begin{proposition}[orders divisible by $2$, $3$ or $5$;\\
  \leanref{A1.even_order_reduce}{AUDIT3/A1_reduction/SizeArith.lean:37},\\
  \leanref{A1.mult3_order_reduce}{AUDIT3/A1_reduction/SizeArith.lean:42},\\
  \leanref{A1.mult5_order_reduce}{AUDIT3/A1_reduction/SizeArith.lean:47}]\label{sa:prop:order}
All variables range over $\N$.
\begin{enumerate}[label=(\alph*)]
\item \textup{(}\lean{A1.even_order_reduce}\textup{)} For all $t,k\in\N$: if $33\le(2t)k$, then $17\le tk$.
\item \textup{(}\lean{A1.mult3_order_reduce}\textup{)} For all $t,k,e\in\N$: if $e\le1$ and $33\le(3t)k+e$, then
  $11\le tk$.
\item \textup{(}\lean{A1.mult5_order_reduce}\textup{)} For all $t,k,e\in\N$: if $e\le1$ and $33\le(5t)k+e$, then
  $35\le5(tk)+e$.
\end{enumerate}
The Lean terms are \mbox{\texttt{2 * t * k}}, \mbox{\texttt{3 * t * k + e}} and \mbox{\texttt{5 * t * k + e}} in the
hypotheses, \mbox{\texttt{t * k}} in the conclusions of \textup{(a)} and \textup{(b)}, and
\mbox{\texttt{5 * (t * k) + e}} in the conclusion of \textup{(c)}.
\end{proposition}

\begin{proof}
By associativity, $(2t)k=2(tk)$, $(3t)k=3(tk)$ and $(5t)k=5(tk)$. Put $m=tk$. Then (a) is
Proposition~\ref{sa:prop:rot}(e) for $m$ (its conclusion $17\le m$), (b) is Proposition~\ref{sa:prop:rot}(a) for $m$,
and (c) is the first conclusion of Proposition~\ref{sa:prop:rot}(b) for $m$. The Lean proofs are the same: each first
proves the associativity identity with \lean{Nat.mul_assoc} and then calls \texttt{omega}, which treats the product
$tk$ as a single unknown.
\end{proof}

\begin{lemma}[the order split;\\
  \leanref{A1.order_split}{AUDIT3/A1_reduction/SizeArith.lean:52}]\label{sa:lem:split}\raggedright
For every $s\in\N$ with $2\le s$: ${s\bmod2=0}$, or ${s\bmod3=0}$, or ${s\bmod5=0}$, or ${7\le s}$.
\end{lemma}

Here $s\bmod q$ is the remainder of $s$ on division by $q$ (Lean \mbox{\texttt{s \% q}}); the Lean conclusion is the
disjunction \mbox{\texttt{s \% 2 = 0}} $\lor$ \mbox{\texttt{s \% 3 = 0}} $\lor$ \mbox{\texttt{s \% 5 = 0}} $\lor$
\mbox{\texttt{7} $\le$ \texttt{s}}.

\begin{proof}
Suppose that $s$ is divisible by none of $2$, $3$, $5$. Then $s\notin\{2,4,6\}$ (these are divisible by $2$),
$s\ne3$ and $s\ne5$. Since $s\ge2$, this leaves $s\ge7$. Lean argues in the same way: it splits on $s<7$ or $s\ge7$
(\lean{Nat.lt_or_ge}); in the first case \texttt{omega} gives $s\in\{2,3,4,5,6\}$, and each of the five resulting
closed instances of the statement is closed by \texttt{decide}, that is, by kernel evaluation of the decision procedure
for the whole closed disjunction (for $s=5$, for example, this evaluates the remainders of $5$ on division by $2$, $3$
and~$5$); in the second case the last alternative holds. The five closed instances are side facts inside the proof
of a statement about all $s\ge2$, which is not itself a finite computation, and the written proof establishes each of
them by hand, through the disjuncts $2\mid2$, $3\mid3$, $2\mid4$, $5\mid5$ and $2\mid6$ used above; hence the lemma
is Lean-checked and is not computer-assisted (see the paragraph ``Rigor'' at the beginning of this section).
\end{proof}

%----------------------------------------------------------------------------------------------------------------------
\subsection{Upper bounds of the size table}\label{sa:ssec:table}

The second part of \dirn{SizeArith.lean} (``Upper bounds of the table, given the case theorems'', lines 58--80) turns
bounds on the number of orbits or pairs into bounds on the size. The case theorems that give such bounds for sets
without $7$ points in convex position are not part of the statements: each part of Proposition~\ref{sa:prop:table}
that rests on a case theorem assumes the bound that the theorem gives as a hypothesis, on $k$ in parts (a), (e)
and~(h), on $2k$ in parts (f) and~(g), on $5k+e$ in part~(c), and on $k$, separately for $b=0$ and for $1\le b$, in
part~(i) (Remark~\ref{sa:rem:use}, item~(5)). Parts (b) and~(d), \lean{A1.c3_no32} and \lean{A1.c5_gap}, rest on no
case theorem: their only hypothesis is $e\le1$.

\begin{proposition}[upper bounds;\\
  \leanref{A1.table_c3}{AUDIT3/A1_reduction/SizeArith.lean:61},\\
  \leanref{A1.c3_no32}{AUDIT3/A1_reduction/SizeArith.lean:62},\\
  \leanref{A1.table_c5}{AUDIT3/A1_reduction/SizeArith.lean:65},\\
  \leanref{A1.c5_gap}{AUDIT3/A1_reduction/SizeArith.lean:66},\\
  \leanref{A1.table_c4}{AUDIT3/A1_reduction/SizeArith.lean:69},\\
  \leanref{A1.table_c4_only_c2}{AUDIT3/A1_reduction/SizeArith.lean:70},\\
  \leanref{A1.table_c6}{AUDIT3/A1_reduction/SizeArith.lean:73},\\
  \leanref{A1.table_c2}{AUDIT3/A1_reduction/SizeArith.lean:76},\\
  \leanref{A1.table_mirror}{AUDIT3/A1_reduction/SizeArith.lean:79}]\label{sa:prop:table}
All variables range over $\N$.
\begin{enumerate}[label=(\alph*)]
\item \textup{(}\lean{A1.table_c3}\textup{)} For all $k,e\in\N$: if $e\le1$ and $k\le10$, then $3k+e\le31$.
\item \textup{(}\lean{A1.c3_no32}\textup{)} For all $k,e\in\N$: if $e\le1$, then $3k+e\ne32$.
\item \textup{(}\lean{A1.table_c5}\textup{)} For all $k,e\in\N$: if $e\le1$ and $5k+e<35$, then $5k+e\le31$.
\item \textup{(}\lean{A1.c5_gap}\textup{)} For all $k,e\in\N$: if $e\le1$, then $5k+e\ne32$, $5k+e\ne33$ and
  $5k+e\ne34$.
\item \textup{(}\lean{A1.table_c4}\textup{)} For every $k\in\N$: if $k\le7$, then $4k\le28$.
\item \textup{(}\lean{A1.table_c4_only_c2}\textup{)} For every $k\in\N$: if $2k\le16$, then $4k\le32$.
\item \textup{(}\lean{A1.table_c6}\textup{)} For every $k\in\N$: if $2k\le10$, then $6k\le30$.
\item \textup{(}\lean{A1.table_c2}\textup{)} For every $k\in\N$: if $k\le16$, then $2k\le32$.
\item \textup{(}\lean{A1.table_mirror}\textup{)} For all $k,b\in\N$: if $b\le2$, $(b=0\Rightarrow k\le16)$ and
  $(1\le b\Rightarrow k\le15)$, then $2k+b\le32$.
\end{enumerate}
\end{proposition}

\begin{proof}
(a) $3k+e\le3\cdot10+1=31$.
(b) If $3k+e=32$, then $3k=32-e\in\{31,32\}$; but $31=3\cdot10+1$ and $32=3\cdot10+2$ are not multiples of $3$.
(c) From $5k\le5k+e<35$ we get $k\le6$, so $5k+e\le30+1=31$.
(d) If $5k+e=j$ with $j\in\{32,33,34\}$, then $5k=j-e$ lies between $31$ and $34$, strictly between the consecutive
multiples $30$ and $35$ of $5$; this is impossible.
(e) $4k\le4\cdot7=28$.
(f) $4k=2\cdot(2k)\le2\cdot16=32$.
(g) $6k=3\cdot(2k)\le3\cdot10=30$.
(h) $2k\le2\cdot16=32$.
(i) If $b=0$, then $k\le16$ and $2k+b=2k\le32$. Otherwise $1\le b$, so $k\le15$, and with $b\le2$ we get
$2k+b\le30+2=32$.
Lean proves each part with \texttt{omega}.
\end{proof}

%----------------------------------------------------------------------------------------------------------------------
\subsection{Size bookkeeping for the assembled statement (audit finding A1-F1)}\label{sa:ssec:assembled}

The file \dirn{AssembledSizes.lean} restates six of the statements above, or parts of them, with the size $n$ as an
explicit variable, adds \lean{big_order_orbit} and a negative control. Its declarations are in the root namespace,
and the file does not import \dirn{SizeArith.lean}.

\begin{proposition}[assembled sizes;\\
  \leanref{halfturn_pairs}{AUDIT3/A1_reduction/verify_A1-F1/AssembledSizes.lean:6},\\
  \leanref{threefold_orbits}{AUDIT3/A1_reduction/verify_A1-F1/AssembledSizes.lean:9},\\
  \leanref{fivefold_size}{AUDIT3/A1_reduction/verify_A1-F1/AssembledSizes.lean:12},\\
  \leanref{mirror_b0}{AUDIT3/A1_reduction/verify_A1-F1/AssembledSizes.lean:15},\\
  \leanref{mirror_b12}{AUDIT3/A1_reduction/verify_A1-F1/AssembledSizes.lean:18},\\
  \leanref{fourfold_via_halfturn}{AUDIT3/A1_reduction/verify_A1-F1/AssembledSizes.lean:21}]\label{sa:prop:assembled}
All variables range over $\N$.
\begin{enumerate}[label=(\alph*)]
\item \textup{(}\lean{halfturn_pairs}\textup{)} For all $n,k\in\N$: if $n=2k$ and $33\le n$, then $17\le k$ and
  $34\le n$.
\item \textup{(}\lean{threefold_orbits}\textup{)} For all $n,k,e\in\N$: if $e\le1$, $n=3k+e$ and $33\le n$, then
  $11\le k$.
\item \textup{(}\lean{fivefold_size}\textup{)} For all $n,k,e\in\N$: if $e\le1$, $n=5k+e$ and $33\le n$, then
  $35\le n$.
\item \textup{(}\lean{mirror_b0}\textup{)} For all $n,k\in\N$: if $n=2k$ and $33\le n$, then $17\le k$.
\item \textup{(}\lean{mirror_b12}\textup{)} For all $n,k,b\in\N$: if $1\le b$, $b\le2$, $n=2k+b$ and $33\le n$, then
  $16\le k$.
\item \textup{(}\lean{fourfold_via_halfturn}\textup{)} For all $n,k\in\N$: if $n=4k$ and $33\le n$, then
  $17\le2k$.
\end{enumerate}
\end{proposition}

\begin{proof}
Substitute the given value of $n$ into $33\le n$. (a) Proposition~\ref{sa:prop:rot}(e) gives $17\le k$ and
$34\le2k=n$. (b) Proposition~\ref{sa:prop:rot}(a). (c) Proposition~\ref{sa:prop:rot}(b) gives $35\le5k+e=n$.
(d) Proposition~\ref{sa:prop:mirror}(a). (e) Proposition~\ref{sa:prop:mirror}(b); as there, $1\le b$ is not used.
(f) Proposition~\ref{sa:prop:rot}(c) gives $17\le2k$. Lean proves each part directly with \texttt{omega}.
\end{proof}

\begin{lemma}[an orbit for an order of at least $7$;\\
  \leanref{big_order_orbit}{AUDIT3/A1_reduction/verify_A1-F1/AssembledSizes.lean:24}]\label{sa:lem:big-order}\raggedright
For all $n,k,e,p\in\N$: if $7\le p$, $e\le1$, $n=pk+e$ and $33\le n$, then $1\le k$.
\end{lemma}

\begin{proof}
If $k=0$, then $n=p\cdot0+e=e\le1<33$, which contradicts $33\le n$. Hence $k\ge1$. We follow the Lean proof, which
splits on $k$ (the product $pk$ is not linear in the unknowns): for $k=0$, \texttt{simp} turns the hypothesis
$n=p\cdot0+e$ into $n=e$, and \texttt{omega} derives a contradiction from $n=e$, $e\le1$ and $33\le n$; for $k=k'+1$
the goal $1\le k'+1$ is closed by \texttt{omega}.
\end{proof}

\begin{remark}[the hypothesis $7\le p$]\label{sa:rem:big-order}
The proof does not use the hypothesis $7\le p$ (Lean binder \texttt{hp}): the conclusion holds for every $p\in\N$,
given $e\le1$, $n=pk+e$ and $33\le n$. The compile log records this as the warning ``unused variable \texttt{hp}''
at line~24, column~41. The name of the lemma and the comment above it (``prime order $p\ge7$'') describe the intended
use (Remark~\ref{sa:rem:use}, item~(1)); the statement does not require $p$ to be prime.
\end{remark}

\begin{remark}[the negative control, \dirn{AssembledSizes.lean:30}]\label{sa:rem:control}
Line~30 of \dirn{AssembledSizes.lean} is an anonymous \texttt{example}, not a named declaration. It states
\[
\neg\bigl(\forall n,k\in\N:\ n=2k\ \Rightarrow\ (32\le n\ \Rightarrow\ 17\le k)\bigr),
\]
that is, Proposition~\ref{sa:prop:assembled}(a) fails when $33$ is replaced by $32$. \emph{Proof.} The numbers $n=32$
and $k=16$ satisfy $32=2\cdot16$ and $32\le32$, but not $17\le16$. The Lean proof is the same:
\texttt{intro h; have := h 32 16 rfl (by decide); omega}, where \texttt{rfl} proves $32=2\cdot16$ by evaluation,
\texttt{decide} proves $32\le32$, and \texttt{omega} refutes $17\le16$. As the example has no name, no
\texttt{\#print axioms} line refers to it; the compile log contains no error, so Lean accepted it. The comment above
the example (line~29) reads ``control: the half-turn line is false below 33 (32 points = 16 pairs; the central
32-point example exists)''. The A1 report describes that example, \dirn{SYM2/sym7_32points.txt}, as half-turn
symmetric about $(0,0)$, not containing the centre, with no collinear triple, and with largest subset in convex
position of $6$ points (\dirn{AUDIT3/reports/A1_reduction.md}, lines 93 and~107); the set is not examined in this
section.
\end{remark}

%----------------------------------------------------------------------------------------------------------------------
\subsection{Baek's theorem at \texorpdfstring{$7$}{7}-gons}\label{sa:ssec:baek}

We recall the notions of Theorem~\ref{mb:thm:baek}. They are defined in Lean in
\dirn{ErdosSzekeres/Vendor/ETV/Config/Defs.lean}, lines 20--67, with the list predicates \lean{List.In} (line~22) and
\texttt{List.Chain3\textquotesingle} (line~46, through the inductive predicate \lean{List.Chain3} of line~42) of
\dirn{ErdosSzekeres/Vendor/ETV/Lib/List/Defs.lean} and the predicate \lean{List.IsChain}, which Mathlib takes from
its dependency Batteries. Let $(\gamma,<)$ be a linearly ordered set. A
\emph{configuration} $C$ on $\gamma$ is a ternary relation $\mathrm{Cup3}_C$ on $\gamma$ (in Lean a structure
\texttt{Config}~$\gamma$ with this relation as a field and a decidability witness, which plays no mathematical role);
$\mathrm{Cap3}_C(a,b,c)$ means $\neg\,\mathrm{Cup3}_C(a,b,c)$. A list $l=(l_0,\dots,l_{r-1})$ of elements of $\gamma$
is a \emph{$C$-cap} if $l_0<l_1<\dots<l_{r-1}$ and $\mathrm{Cap3}_C(l_i,l_{i+1},l_{i+2})$ for every $i$ with $i+2<r$,
and a \emph{$C$-cup} if the same holds with $\mathrm{Cup3}_C$ in place of $\mathrm{Cap3}_C$. For a finite set
$S\subseteq\gamma$, $\mathsf{HasNCap}_C(4,S)$ means that some $C$-cap of length $4$ has all its entries in $S$, and
$\mathsf{HasNGon}_C(7,S)$ means that there are lists $l$ and $l'$ with all entries in $S$ such that $\card l\ge2$, $l$
is a $C$-cap, $\card{l'}\ge2$, $l'$ is a $C$-cup, $l$ and $l'$ have the same first entry and the same last entry, and
$\card l+\card{l'}=7+2$. The relation $\mathrm{Cup3}_C$ is arbitrary, so these notions carry no geometry.

\begin{theorem}[Baek's theorem at $7$-gons;\\
  \leanref{baek_seven}{ES7SPLIT/review/equivalence/BaekSeven.lean:6}]\label{sa:thm:baek-seven}
Let $\gamma$ be a type \textup{(}Lean \mbox{\texttt{$\gamma$ : Type}}, the lowest universe\textup{)} with a linear
order $<$, let $C$ be a configuration on $\gamma$ \textup{(}\mbox{\texttt{C : Config}~$\gamma$}\textup{)}, and let
$S$ be
a finite subset of $\gamma$ \textup{(}\mbox{\texttt{S : Finset}~$\gamma$}\textup{)} with $17\le\card S$. Then
$\mathsf{HasNCap}_C(4,S)$ or $\mathsf{HasNGon}_C(7,S)$.
\end{theorem}

The theorem is Lean-checked (\dirn{BaekSeven.out}, line~3). Its written proof rests on the cited
Theorem~\ref{mb:thm:baek}.

\begin{proof}
We have $\binom{4+2}{2}+2=\binom62+2=\frac{6\cdot5}{2}+2=15+2=17\le\card S$. Theorem~\ref{mb:thm:baek} with $n=4$
and the same $\gamma$, $C$ and $S$ gives $\mathsf{HasNCap}_C(4,S)$ or $\mathsf{HasNGon}_C(4+3,S)$, and $4+3=7$. The
Lean proof
is the same. Its first line, \texttt{have e : Nat.choose (4 + 2) 2 + 2 = 17 := by decide}, evaluates the binomial
coefficient in the kernel (the value $\binom62=15$ computed above by hand). Its second line,
\texttt{exact main 4 C S (e} $\blacktriangleright$ \texttt{h)}, applies \lean{main} with $n=4$: the term
\texttt{e} $\blacktriangleright$ \texttt{h} (Lean's rewriting operator) turns the hypothesis \texttt{h}, which says
$17\le\card S$, into the hypothesis $\binom{4+2}{2}+2\le\card S$ of \lean{main} at $n=4$, and Lean accepts the
conclusion \mbox{\texttt{C.HasNGon (4 + 3) S}} as \mbox{\texttt{C.HasNGon 7 S}} because $4+3$ evaluates to $7$.
\end{proof}

\begin{remark}[comparison with Theorem~\ref{mb:thm:baek}; scope]\label{sa:rem:baek-seven}
(1) Lines 1--2 of \dirn{BaekSeven.out} are the output of \mbox{\texttt{\#check @main}}. They say that \lean{main}
takes, in this order, an implicit type $\gamma$ in an arbitrary universe (\mbox{\texttt{Type u\_1}}), an instance
\mbox{\texttt{LinearOrder}~$\gamma$}, $n\in\N$, \mbox{\texttt{C : Config}~$\gamma$} and
\mbox{\texttt{S : Finset}~$\gamma$}, and that it derives from the hypothesis
\mbox{\texttt{(n + 2).choose 2 + 2} $\le$ \texttt{S.card}}, that is $\binom{n+2}{2}+2\le\card S$, the conclusion
\mbox{\texttt{C.HasNCap 4 S}} $\lor$ \mbox{\texttt{C.HasNGon (n + 3) S}}. This is the statement of
Theorem~\ref{mb:thm:baek}, with the same notions; it is also the declaration at
\dirn{ErdosSzekeres/Vendor/ETV/Main/Main.lean:87} together with the variables declared at line~26 of that file. The
cap length $4$ is fixed; only the gon size $n+3$ and the threshold depend on $n$.

(2) Theorem~\ref{sa:thm:baek-seven} is the case $n=4$. The only other difference is the universe: \lean{baek_seven}
is stated for $\gamma$ in the lowest universe \texttt{Type}, \lean{main} for types in every universe. Read
mathematically, both statements are about an arbitrary linearly ordered set.

(3) The theorem concerns abstract configurations; it mentions no points of the plane, and
\dirn{BaekSeven.lean} contains no bridge from point sets to configurations. The records of \dirn{ES7SPLIT} say:
``The step from the abstract form to real point sets is not Lean-checked for $n = 7$; the repository's bridge covers
the octagon case only. Here that step rests on the published paper.'' (\dirn{ES7SPLIT/README.md}, lines 78--79, in
the rigor paragraph of its section~2; the status paragraph of \dirn{ES7SPLIT/proofs/split_theorem.md}, lines 13--15,
says the same.) These records were committed in \dirn{9793c06e0} (2~November 2026, 01:03 +0900) and predate
\lean{MirrorOdd.capacity16} (\dirn{MIRRORETV/lean/MirrorMain.lean:1529}, added in \dirn{83c80f38f} at 03:41 the same
day), a geometric form of the same instance of Theorem~\ref{mb:thm:baek}: a configuration of $n$ points in the sense
of Definition~\ref{mm:def:config} with no $4$-cap and no $7$ points in convex position has $n\le16$. That theorem is
Lean-checked: the Lean proof of Theorem~\ref{mb:thm:seven} uses it, and the axioms printed for that theorem cover it.
Whether it covers the use of Baek's theorem in the $z$-split arguments of \dirn{ES7SPLIT} is not examined here; the
theorems of \dirn{ES7SPLIT} remain written proofs. The source of Baek's theorem is J.~Baek, \emph{On the
Erd\H{o}s--Tuza--Valtr conjecture}, Eur.\ J.\ Combin.\ 124 (2025) 104085; arXiv:2206.04260.
\end{remark}

%----------------------------------------------------------------------------------------------------------------------
\subsection{How the arithmetic is used}\label{sa:ssec:use}

The following lemma is the modelling of the header comment of \dirn{SizeArith.lean}, proved here as a written step.
For $c\in\R^2$ and $\theta\in\R$, $\rho_{c,\theta}$ is the rotation of Definition~\ref{fd:def:rot}; in complex
notation it is $z\mapsto c+e^{i\theta}(z-c)$. For a map $f$ we write $f^{[i]}$ for its $i$-th iterate. The names
\lean{HalfTurn34}, \lean{Threefold33}, \lean{Fivefold35}, \lean{MirrorOneAxis33}, \lean{MirrorNoAxis34} and
\lean{Fourfold32} denote the case classes of Definition~\ref{sr:def:classes}, in the namespace
\lean{ErdosSzekeres.SymReduction}.

\begin{lemma}[orbit counts; written step]\label{sa:lem:model}
Let $S\subseteq\R^2$ be a finite set and $n=\card S$.
\begin{enumerate}[label=(\alph*)]
\item \emph{Rotations.} Let $c\in\R^2$, let $s\in\N$ with $0<s$, let $\rho=\rho_{c,2\pi/s}$, and assume
  $\mathrm{Invariant}\,\rho\,S$ \textup{(}Definition~\ref{fd:def:inv}\textup{)}. For $y\in\R^2$ put
  $O(y)=\{\rho^{[i]}(y):0\le i<s\}$. For every $y\in S\setminus\{c\}$ the set $O(y)$ is contained in
  $S\setminus\{c\}$, has exactly $s$ elements and is mapped into itself by $\rho$; any two of these sets are equal or
  disjoint; and $n=s\cdot k+e$, where $k$ is the number of distinct sets $O(y)$ with $y\in S\setminus\{c\}$, and $e=1$
  if $c\in S$, $e=0$ otherwise. If moreover $s$ is even, $S$ is in general position and $S\not\subseteq\{c\}$, then
  $e=0$.
\item \emph{Reflections.} Let $a,u\in\R^2$ with $u_0^2+u_1^2=1$, let $\sigma=\mathrm{reflLine}_{a,u}$, that is
  \[
  \sigma(p)=\bigl(a_0+2du_0-(p_0-a_0),\;a_1+2du_1-(p_1-a_1)\bigr),\qquad d=(p_0-a_0)u_0+(p_1-a_1)u_1,
  \]
  the map of the classes \lean{MirrorOneAxis33} and \lean{MirrorNoAxis34} of Definition~\ref{sr:def:classes}, and
  let $L=\{a+tu:t\in\R\}$. Assume $\mathrm{Invariant}\,\sigma\,S$ and that $S$ is in general position. Then
  $S\setminus L$ is the disjoint union of $k$ two-element sets of the form $\{p,\sigma(p)\}$, each mapped into itself
  by $\sigma$, and $n=2k+b$ with $b=\card{S\cap L}\le2$.
\item \emph{Deleting.} Let $T\subseteq S$. If $S$ is in general position, so is $T$; and
  $\mathrm{HasConvexNGon}\,7\,T$ implies $\mathrm{HasConvexNGon}\,7\,S$. In \textup{(a)}, a union of sets $O(y)$,
  together with $c$ or not when $c\in S$, is mapped into itself by $\rho$; in \textup{(b)}, a union of sets
  $\{p,\sigma(p)\}$ and of points of $S\cap L$ is mapped into itself by $\sigma$.
\end{enumerate}
\end{lemma}

\begin{proof}
(a) Put $\omega=e^{2\pi i/s}$, so that $\rho(z)=c+\omega(z-c)$. Then $\rho(c)=c$, and $\rho$ is injective because
$\omega\ne0$; so every iterate $\rho^{[i]}$ is injective. Hence $\rho$ maps $S\setminus\{c\}$ into itself: for
$y\in S\setminus\{c\}$ we have $\rho(y)\in S$, and $\rho(y)=c=\rho(c)$ would give $y=c$. By induction on $i$,
$\rho^{[i]}(y)\in S\setminus\{c\}$ for all $i$, so $O(y)\subseteq S\setminus\{c\}$. Let $y\in S\setminus\{c\}$. By
Lemma~\ref{ll:lem:minperiod} (with $m=s$), the minimal period of $y$ under $\rho$ is $s$: $\rho^{[s]}(y)=y$ and
$\rho^{[i]}(y)\ne y$ for $0<i<s$. The points $\rho^{[i]}(y)$, $0\le i<s$, are pairwise distinct: if
$\rho^{[i]}(y)=\rho^{[j]}(y)$ with $0\le i<j<s$, then $\rho^{[i]}\bigl(\rho^{[j-i]}(y)\bigr)=\rho^{[i]}(y)$, and
injectivity of $\rho^{[i]}$ gives $\rho^{[j-i]}(y)=y$ with $0<j-i<s$, a contradiction. So $\card{O(y)}=s$. Further,
$\rho(\rho^{[i]}(y))=\rho^{[i+1]}(y)\in O(y)$ for $i+1<s$, and $\rho(\rho^{[s-1]}(y))=\rho^{[s]}(y)=y\in O(y)$; so
$\rho$ maps $O(y)$ into itself. Consequently, for $x\in O(y)$ all points $\rho^{[i]}(x)$ lie in $O(y)$, so
$O(x)\subseteq O(y)$, and $O(x)=O(y)$ because both sets have $s$ elements. If $z\in O(y)\cap O(y')$, then
$O(y)=O(z)=O(y')$. Every $y\in S\setminus\{c\}$ lies in $O(y)$ (take $i=0$). Therefore the distinct sets $O(y)$ form a
partition of $S\setminus\{c\}$ into $k$ blocks of size $s$, and $n=\card{S\setminus\{c\}}+e=s\cdot k+e$. If $s$ is
even, $S$ is in general position and some $p\in S$ differs from $c$, then Lemma~\ref{lg:lem:centre} (with $m=s$)
gives $c\notin S$, that is, $e=0$.

(b) Write $\langle x,y\rangle=x_0y_0+x_1y_1$, so that $d=\langle p-a,u\rangle$ and $\sigma(p)-a=2du-(p-a)$.
\emph{$\sigma$ is an involution:} for $p'=\sigma(p)$ we get $\langle p'-a,u\rangle=2d\langle u,u\rangle-d=d$, hence
$\sigma(p')-a=2du-(p'-a)=2du-2du+(p-a)=p-a$, that is, $\sigma(\sigma(p))=p$; in particular $\sigma$ is injective.
\emph{The fixed points of $\sigma$ are the points of $L$:} $\sigma(p)=p$ means $2du=2(p-a)$, so $p=a+du\in L$;
conversely, for $p=a+tu$ we have $d=t\langle u,u\rangle=t$ and $\sigma(p)=a+2tu-tu=p$. (This is also shown at the
end of Definition~\ref{sr:def:classes}.) Now let $p\in S\setminus L$. Then $\sigma(p)\in S$ and $\sigma(p)\ne p$; and
$\sigma(p)\notin L$, since otherwise $\sigma(\sigma(p))=\sigma(p)$, that is, $p=\sigma(p)$. So $\{p,\sigma(p)\}$ is a
two-element subset of $S\setminus L$, mapped into itself by $\sigma$ because $\sigma(\sigma(p))=p$. If two such sets
$\{p,\sigma(p)\}$ and $\{q,\sigma(q)\}$ meet, then $q$ or $\sigma(q)$ lies in $\{p,\sigma(p)\}$, and applying
$\sigma$ shows that both $q$ and $\sigma(q)$ do; so the two sets are equal. Hence the distinct sets $\{p,\sigma(p)\}$
partition $S\setminus L$, $\card{S\setminus L}=2k$, and $n=2k+b$ with $b=\card{S\cap L}$. Finally $b\le2$: if
$x,y,z\in S\cap L$ were pairwise distinct, all differences of two of them would be real multiples of $u$, so the
linear span of these differences would have dimension at most $1$ and $\{x,y,z\}$ would be collinear, contrary to
general position (Definition~\ref{fd:def:gp}). This is the fact that at most two points of a set in general position
lie on a line.

(c) The first two claims are Lemma~\ref{fd:lem:mono}(b) and~(c). A union of sets each of which is mapped into itself
by a map is mapped into itself by that map; $\{c\}$ is mapped into itself by $\rho$ since $\rho(c)=c$, and each point
of $S\cap L$ is fixed by $\sigma$ by (b).
\end{proof}

\begin{remark}[How the arithmetic is used]\label{sa:rem:use}
This remark is a written step; it is not Lean-checked. It follows the reduction re-derived by the auditor in
\dirn{AUDIT3/reports/A1_reduction.md} (Summary and trust-table rows 2--9) and shows where each statement of this
section enters. Let $S$ be a finite set in general position with $n=\card S\ge33$.

(1) \emph{Rotations.} Let $c\in\R^2$, $s\ge2$, $\rho=\rho_{c,2\pi/s}$, and let $S$ be mapped into itself by $\rho$.
By Lemma~\ref{sa:lem:model}(a), $n=sk+e$ with $e\le1$, and $e=0$ for even $s$ (as $n\ge2$, $S\not\subseteq\{c\}$). By
Lemma~\ref{sa:lem:model}(c), a union $T$ of some of the sets $O(y)$ is in general position, mapped into itself by
$\rho$ and does not contain $c$.
\begin{itemize}
\item $s=2$: as $\cos\pi=-1$ and $\sin\pi=0$, $\rho=\rho_{c,\pi}$ is the half-turn $p\mapsto2c-p$ (Lean
  \mbox{\texttt{halfTurn c}}), and $n=2k$. Proposition~\ref{sa:prop:rot}(e) (or~\ref{sa:prop:assembled}(a)) gives
  $k\ge17$ and $n\ge34$. A union $T$ of $17$ pairs has $34$ points: these are the hypotheses of \lean{HalfTurn34}
  (Definition~\ref{sr:def:classes}) for $T$ and $c$.
\item $s=3$: Proposition~\ref{sa:prop:rot}(a) (or~\ref{sa:prop:assembled}(b)) gives $k\ge11$; $11$ orbits give
  $33$ points without $c$, invariant under $\rho_{c,2\pi/3}$: the hypotheses of \lean{Threefold33}.
\item $s=5$: Proposition~\ref{sa:prop:rot}(b) gives $n\ge35$ and $k\ge7$ (Proposition~\ref{sa:prop:assembled}(c)
  gives $n\ge35$ only); $7$ orbits give $35$ points without $c$: the hypotheses of \lean{Fivefold35}. (The comment of
  \lean{A1.c5_reduce}, ``$T_{C5}$ applies directly'', refers to the fivefold case theorem of the records, which admits
  at least $35$ points and the centre; the Lean class asks for exactly $35$ points without the centre, which $k\ge7$
  supplies.)
\item $s=4$: Proposition~\ref{sa:prop:rot}(c) gives $n=4k\ge36$, $k\ge8$ and $2k\ge17$
  (Proposition~\ref{sa:prop:assembled}(f) gives $2k\ge17$ only). Rotations about $c$ compose by adding angles (in
  complex notation, $c+e^{i\alpha}\bigl(c+e^{i\beta}(z-c)-c\bigr)=c+e^{i(\alpha+\beta)}(z-c)$), so
  $\rho^{[2]}=\rho_{c,\pi}$ is the
  half-turn; Lemma~\ref{sa:lem:model}(a) for it gives $n=2k'$, and $k'=2k\ge17$, so the case $s=2$ applies. (Eight of
  the $k$ orbits would give the $32$ points of the class \lean{Fourfold32}. Among the reduction theorems and case
  lemmas stated in this document, only Theorem~\ref{sr:thm:variant} and the case lemma
  \lean{SymReduction.case_fourfold} that its proof uses take this class as a hypothesis (so do some checks of that
  theorem in the review file \dirn{SYMLEAN/review/VariantChecks.lean}, among them its group form
  \lean{SymReduction.VariantChecks.group_form}); a same-text copy, \lean{LayerLemma.Fourfold32}, is the hypothesis of
  Theorem~\ref{ec:thm:fourfold} and of its layer form \lean{LayerLemma.fourfold_heptagon_of_erase}, which its proof
  uses;
  item~(5) uses the class for the size table.)
\item $s=6$: Proposition~\ref{sa:prop:rot}(d) gives $2k\ge11$ and $3k\ge17$. Here $\rho^{[2]}=\rho_{c,2\pi/3}$ has
  $2k$ orbits of size $3$ and $\rho^{[3]}=\rho_{c,\pi}$ has $3k$ pairs (Lemma~\ref{sa:lem:model}(a) for these
  rotations, with $n=6k$), so both the threefold and the half-turn case apply.
\item General $s$: by Lemma~\ref{sa:lem:split}, $2\mid s$, $3\mid s$, $5\mid s$, or $s\ge7$. If $s=mt$ with
  $m\in\{2,3,5\}$, then $\rho^{[t]}=\rho_{c,2\pi t/s}=\rho_{c,2\pi/m}$ maps $S$ into itself, and
  Lemma~\ref{sa:lem:model}(a) for it gives $n=mk'+e$ with the same $e$; from $mk'=n-e=mtk$ we get $k'=tk$.
  Proposition~\ref{sa:prop:order} is the arithmetic of this step: $tk\ge17$ pairs for even $s$ (where $e=0$),
  $tk\ge11$ orbits of size $3$ if $3\mid s$, and $35\le5(tk)+e$ if $5\mid s$, hence $5tk\ge34$ and $tk\ge7$. The cases
  $s=2,3,5$ then apply to $\rho^{[t]}$. If $s\ge7$, Lemma~\ref{sa:lem:big-order} (with $p=s$) gives $k\ge1$, so $S$
  contains a set $O(y)$ with $s\ge7$ points. They all have the same distance from $c$ as $y$, because
  $\abs{\omega(z-c)}=\abs{z-c}$, so $O(y)$ is in convex position by Lemma~\ref{sr:lem:concyclic}; any $7$ of its
  points are in convex position (Lemma~\ref{fd:lem:mono}(a)), so $\mathrm{HasConvexNGon}\,7\,S$
  (Definition~\ref{fd:def:ngon}). No case class is needed here.
\end{itemize}

(2) \emph{Reflections.} Let $S$ be mapped into itself by $\sigma=\mathrm{reflLine}_{a,u}$ with $u_0^2+u_1^2=1$. By
Lemma~\ref{sa:lem:model}(b), $n=2k+b$ with $b\le2$. If $b=0$, Proposition~\ref{sa:prop:mirror}(a)
(or~\ref{sa:prop:assembled}(d)) gives $k\ge17$, and $17$ pairs form a $34$-point set that is in general position and
mapped into itself by $\sigma$ (Lemma~\ref{sa:lem:model}(c)) and has no fixed point: the hypotheses of
\lean{MirrorNoAxis34}. If $b\in\{1,2\}$, Proposition~\ref{sa:prop:mirror}(b) (or~\ref{sa:prop:assembled}(e)) gives
$k\ge16$, and $16$ pairs together with one point $x\in S\cap L$ form a $33$-point set that is in general position and
mapped into itself by $\sigma$ (Lemma~\ref{sa:lem:model}(c)) and whose only fixed point is $x$: the hypotheses of
\lean{MirrorOneAxis33}. (The fixed points of $\sigma$ are the points of $L$, as shown in the proof of
Lemma~\ref{sa:lem:model}(b), and the pairs lie in $S\setminus L$.)

(3) \emph{All symmetries.} By Theorem~\ref{sr:thm:structure}(b), a Euclidean isometry $g\ne\mathrm{id}$ that maps $S$
into itself is either a rotation, in complex notation $z\mapsto c+a(z-c)$ with $a\ne1$ and $a^m=1$ for some $m>0$, or
a map $\mathrm{reflLine}_{A,u}$ with $u_0^2+u_1^2=1$. In the first case let $s\ge2$ be the order of $a$. Then $a$ is a
primitive $s$-th root of unity, so $e^{2\pi i/s}=a^j$ for some $j\in\N$, because a primitive $s$-th root of unity
generates all $s$-th roots of unity. As $g^{[j]}(z)=c+a^j(z-c)$ (induction on $j$), $\rho_{c,2\pi/s}=g^{[j]}$, which
maps $S$ into itself (here $c$ is read as a point); so item~(1) applies. In the second case item~(2) applies.
Affine symmetries are reduced to Euclidean isometries by Theorem~\ref{ar:thm:transfer}, which gives
Theorem~\ref{ar:thm:main}.

(4) \emph{Status.} Items (1)--(3) derive $\mathrm{HasConvexNGon}\,7\,S$ from the five case classes \lean{HalfTurn34},
\lean{Threefold33}, \lean{Fivefold35}, \lean{MirrorOneAxis33} and \lean{MirrorNoAxis34}, used as hypotheses. Except
for the orders $s\ge7$ in item~(1), which need no class, items (1) and~(2) exhibit a subset $T$ that satisfies the
hypotheses of one class, and the class then gives $\mathrm{HasConvexNGon}\,7\,T$, hence
$\mathrm{HasConvexNGon}\,7\,S$ by Lemma~\ref{sa:lem:model}(c); item~(3) reduces every Euclidean isometry other than
the identity to items (1) and~(2), and affine symmetries to Euclidean isometries. The Lean-checked form of this
implication is Theorem~\ref{sr:thm:main} for Euclidean isometries, with Theorem~\ref{ar:thm:main} for affine
symmetries. The Lean proof of Theorem~\ref{sr:thm:main} passes to a power of the rotation that is the half-turn or has
odd prime order, instead of using Lemma~\ref{sa:lem:split}, and it does its own arithmetic. The two mirror classes
are Lean-checked theorems (Theorem~\ref{mb:thm:classes}). The classes \lean{HalfTurn34}, \lean{Threefold33} and
\lean{Fivefold35} are not proved in Lean or in this document (Theorem~\ref{mb:thm:main3} is conditional on them);
the records that supply them are described in the remark that follows Definition~\ref{sr:def:classes}. Of the
three ingredients, the modelling (Lemma~\ref{sa:lem:model}), the case classes and the arithmetic, the files of this
section check only the arithmetic.

(5) \emph{The size table.} Proposition~\ref{sa:prop:table} is the arithmetic of the upper bounds in the size table of
\dirn{WRITEUP/SYMMETRIC.md}, \S2 (trust-table row~9 of the A1 report). For a set without $7$ points in convex position,
the case classes bound $k$ through item~(1) or~(2), and the parts of Proposition~\ref{sa:prop:table} turn these
bounds into bounds on $n$. Here $S$ need not have $33$ points: if the number of orbits or pairs to which a class is
applied below ($k$, or $2k$ for $\rho^{[2]}$ when $s=4$ or $s=6$) reached the number that the class requires, the
subset $T$ built in item~(1) or~(2) would satisfy the hypotheses of the class, which would then give $7$ points of
$S$ in convex position; for even $s$, $e=0$ as soon as $k\ge1$ (Lemma~\ref{sa:lem:model}(a)), and $k=0$ gives
$n=e\le1$. For threefold sets, \lean{Threefold33} gives $k\le10$, so $n\le31$ by part~(a), and
$n\ne32$ in any case by part~(b). For fivefold sets, \lean{Fivefold35} gives $k\le6$, so $n=5k+e<35$ and $n\le31$ by
part~(c); $32$, $33$ and $34$ are not fivefold sizes by part~(d). For fourfold sets, \lean{Fourfold32} gives $k\le7$,
so $n\le28$ by part~(e), while \lean{HalfTurn34} alone gives only $2k\le16$ and $n\le32$ by part~(f). For sixfold
sets, \lean{Threefold33} applied to the $2k$ orbits of size $3$ gives $2k\le10$, so $n\le30$ by part~(g). For
half-turns, \lean{HalfTurn34} gives $k\le16$, so $n\le32$ by part~(h). For reflections, the two mirror classes give
$k\le16$ if $b=0$ and $k\le15$ if $b\ge1$, so $n\le32$ by part~(i).
Except for the mirror row, these bounds are conditional on the classes named, which are hypotheses here: the rows
for half-turns, threefold and sixfold sets and the bound $32$ of part~(f) on \lean{HalfTurn34} or \lean{Threefold33},
the fourfold bound $28$ on \lean{Fourfold32}, and the fivefold bound $31$ on \lean{Fivefold35}. For the last two,
Lean-checked conditional forms exist that include the geometric modelling: Theorem~\ref{ec:thm:fourfold}, with the
same-text copy \lean{LayerLemma.Fourfold32} as hypothesis, gives $\mathrm{HasConvexNGon}\,7\,S$ for every set $S$ in
general position with at least $29$ points that is mapped into itself by $\rho_{c,2\pi/4}$, so such a set without $7$
points in convex position has at most $28$ points; Theorem~\ref{ec:thm:fivefold}, with the same-text copy
\lean{LayerLemma.Fivefold35}, does the same with $\rho_{c,2\pi/5}$, at least $32$ points and the bound $31$. For
reflections the bound $n\le32$ is Lean-checked with no class hypothesis (Theorem~\ref{mb:thm:refl32}(b) for
$\mathrm{reflLine}_{a,u}$, derived from Theorem~\ref{mb:thm:seven}(b)), and the two mirror classes are themselves
Lean-checked theorems (Theorem~\ref{mb:thm:classes}). The matching lower bounds rest on example sets of the records.
None of this is part of the files of this section.
\end{remark}

{\raggedright
\paragraph{Coverage.}
Covered, by label. In \dirn{AUDIT3/A1_reduction/SizeArith.lean}: \lean{A1.c3_reduce}, \lean{A1.c5_reduce},
\lean{A1.c4_reduce}, \lean{A1.c6_reduce}, \lean{A1.c2_reduce} (Proposition~\ref{sa:prop:rot});
\lean{A1.mirror_reduce_b0}, \lean{A1.mirror_reduce_b12} (Proposition~\ref{sa:prop:mirror});
\lean{A1.even_order_reduce}, \lean{A1.mult3_order_reduce}, \lean{A1.mult5_order_reduce}
(Proposition~\ref{sa:prop:order}); \lean{A1.order_split} (Lemma~\ref{sa:lem:split}); \lean{A1.table_c3},
\lean{A1.c3_no32}, \lean{A1.table_c5}, \lean{A1.c5_gap}, \lean{A1.table_c4}, \lean{A1.table_c4_only_c2},
\lean{A1.table_c6}, \lean{A1.table_c2}, \lean{A1.table_mirror} (Proposition~\ref{sa:prop:table}). In
\dirn{AUDIT3/A1_reduction/verify_A1-F1/AssembledSizes.lean}: \lean{halfturn_pairs}, \lean{threefold_orbits},
\lean{fivefold_size}, \lean{mirror_b0}, \lean{mirror_b12}, \lean{fourfold_via_halfturn}
(Proposition~\ref{sa:prop:assembled}); \lean{big_order_orbit} (Lemma~\ref{sa:lem:big-order}, with
Remark~\ref{sa:rem:big-order}); the anonymous example of line~30 (Remark~\ref{sa:rem:control}). In
\dirn{ES7SPLIT/review/equivalence/BaekSeven.lean}: \lean{baek_seven} (Theorem~\ref{sa:thm:baek-seven}); the vendored
theorem \lean{main} that it instantiates is the cited Theorem~\ref{mb:thm:baek}. That is all $28$ theorems of the
three files and the one example.

Not covered, deliberately: the \texttt{\#print axioms} commands (\dirn{SizeArith.lean}, lines 84--103;
\dirn{AssembledSizes.lean}, lines 33--39; \dirn{BaekSeven.lean}, line~12) and the \texttt{\#check} command
(\dirn{BaekSeven.lean}, line~11), which produce the logs named in the Sources paragraph and are not results; and the
comments of the files, among them the modelling in the header of \dirn{SizeArith.lean}, which these files do not
formalise and which is treated here as the written step Lemma~\ref{sa:lem:model}. Lemma~\ref{sa:lem:model} and
Remark~\ref{sa:rem:use} are not declarations of the three files. Out of scope: the scripts and other logs of
\dirn{AUDIT3/A1_reduction} and its subdirectories, and the other files of \dirn{ES7SPLIT/review/equivalence}; none of
them contains a Lean declaration.\par}
